\begin{lemma}[Noncentral vacuum density from an untwisted boundary component]
    \label{lem:peps_anchored_twist_density}
    \lean{TNLean.PEPS.regularPhysicalCutMatrix_normalized_reduced_eq_of_untwisted_boundary_component}
    \leanok
    Let the physical sites be regular group-isometric. Suppose the observed
    and complementary induced graphs are connected, the boundary is
    nonempty, and every insertion on an observed incident bond is the
    identity. If one untwisted complementary component contains every
    complementary boundary endpoint, both actual cuts are nonzero and
    their normalized observed reduced densities agree:
    \begin{align}
        \frac{MM^\dagger}{\operatorname{tr}(MM^\dagger)}
        =\frac{M_0M_0^\dagger}{\operatorname{tr}(M_0M_0^\dagger)}.
        \label{eq:peps_anchored_twist_density}
    \end{align}
    Here $M$ is the inserted cut and $M_0$ the vacuum cut.
    No centrality of the group insertions is required.
\end{lemma}
\begin{proof}\leanok
    \uses{lem:peps_anchored_complement_gauge,thm:peps_regular_twisted_state_nonzero,
        thm:peps_regular_twisted_cut_entropy,thm:peps_regular_boundary_normalization,
        thm:peps_regular_boundary_physical_realization,thm:peps_fixed_boundary_factor_density}
    Construct the complementary tree gauge, extend it by the identity on
    the observation region, and preserve the actual physical vector.
    The resulting complementary tree is a common connected untwisted
    subgraph for the transformed insertion and the vacuum. The projected
    mixed Gram calculation from the local isometries gives a common
    invariant-boundary density $\rho$ with $\operatorname{tr}\rho=1$.
    For nonzero scalars $c,c_0$ the calculation gives
    $MM^\dagger=c\rho$ and $M_0M_0^\dagger=c_0\rho$.
    Thus $c=\operatorname{tr}(MM^\dagger)$ and
    $c_0=\operatorname{tr}(M_0M_0^\dagger)$, and
    \begin{align}
        \operatorname{tr}(MM^\dagger)^{-1}MM^\dagger
        =\rho
        =\operatorname{tr}(M_0M_0^\dagger)^{-1}M_0M_0^\dagger.
        \label{eq:peps_anchored_twist_common_density}
    \end{align} Local physical
    inverses expose a nonzero canonical projector column, proving
    nonvanishing of both actual cuts.
\end{proof}
